% META: {"id": "TSPE-PRIMEQ-CR-013", "chapitre": "TSPE-PRIMITIVES-EQDIFF", "type_objet": "cours", "capacites_codes": ["C2"], "sources_inspiration": [], "mode_creation": "ex_nihilo", "status": "approved"}

\section{Équation différentielle $y' = ay + b$}

% ============================================================
% STRATE 1 — Méthode et théorème
% ============================================================

\theoreme{
Soit $a \in \mathbb{R}^*$ et $b \in \mathbb{R}$. L'équation $y' = ay+b$ admet une unique solution constante $y_0 = -\dfrac{b}{a}$ (solution particulière).

Les solutions de $y'=ay+b$ sur $\mathbb{R}$ sont exactement les fonctions :
\[
  x \mapsto C\,\mathrm{e}^{ax} - \frac{b}{a}, \qquad C \in \mathbb{R}.
\]
}

\margeAppui{
\textbf{Méthode.} On cherche d'abord la solution particulière constante $y_0$ (verifiant $y_0'=0=ay_0+b$, d'ou $y_0=-b/a$). Puis on pose $z = y-y_0$ : $z' = y' = ay+b = a(z+y_0)+b = az + (ay_0+b) = az$ (car $ay_0+b=0$). Donc $z$ vérifie $z'=az$, et $z(x)=C\mathrm{e}^{ax}$ ; d'ou $y = z+y_0 = C\mathrm{e}^{ax} - b/a$.
}

\exemple{
Résoudre $y' = 2y - 6$, puis déterminer la solution telle que $y(0)=1$.

Solution particulière constante : $y_0 = -\dfrac{-6}{2} = 3$ (on vérifie : $y_0'=0$ et $2y_0-6=6-6=0$. \checkmark)

Solutions generales : $y(x) = C\mathrm{e}^{2x} + 3$.

Condition $y(0)=1$ : $C\mathrm{e}^0+3 = C+3 = 1$, donc $C=-2$.

La solution cherchee est $y(x) = -2\mathrm{e}^{2x}+3$.

% BEGIN-VERIFY
% from sympy import symbols, exp, diff
% x = symbols('x')
% y = -2*exp(2*x) + 3
% assert diff(y, x) == 2*y - 6
% assert y.subs(x, 0) == 1
% END-VERIFY
}

% ============================================================
% STRATE 2 — Erreurs frequentes
% ============================================================

\erreurFrequente{
\textbf{Se tromper de signe pour la solution particulière.} Pour $y'=ay+b$, la solution constante est $y_0 = -\dfrac{b}{a}$ (avec un signe moins), obtenue en resolvant $0 = ay_0+b$. Une erreur frequente est d'écrire $y_0 = \dfrac{b}{a}$.
}
